\section{Familias predominantes de Actor-Critic y diagnóstico}

\subsection{Familias de algoritmos y diferencias de diseño}

\begin{table}[H]
\centering
\begin{tabular}{p{0.23\textwidth}p{0.21\textwidth}p{0.23\textwidth}p{0.2\textwidth}}
\toprule
Algoritmo & Tipo de Critic & Actualización de la política & Escenario adecuado \\
\midrule
PPO & soft target de V & clipped surrogate & entrenamiento síncrono de modelos, clústeres de CPU \\
A2C/A3C & V de un paso + entropy & sobremuestreo con distintos workers & paralelismo coordinado, múltiples instancias \\
SAC & doble Q + entropy max & off-policy replay buffer & acciones continuas, que requieren alta eficiencia de muestra \\
DDPG/TD3 & Deterministic policy & deterministic actor gradient & control continuo de alta dimensión \\
\bottomrule
\end{tabular}
\caption{Diferencias centrales entre las distintas variantes de Actor-Critic}
\end{table}

\begin{knowledgebox}{Dimensiones de diseño abstraídas de la diapositiva de la clase}
Las diferencias que el docente enfatiza provienen principalmente de tres dimensiones: la frecuencia de actualización (on-policy vs off-policy), la restricción de la función objetivo (clipping / KL / trust region), y si incluye entropy o ajusta automáticamente la exploration.
\end{knowledgebox}

\begin{figure}[H]
\centering
\includegraphics[width=0.85\textwidth]{slide-12.jpg}
\caption{Figura comparativa de PPO y SAC de las Lecture slides, que enfatiza la complementariedad entre la estabilidad de on-policy y la eficiencia de muestra de off-policy.}
\end{figure}

\subsection{Monitoreo del entrenamiento y métricas experimentales}

\begin{itemize}
    \item \textbf{KL divergence}: monitorea si la actualización rebasa el trust region;
    \item \textbf{Distribución de Advantage}: si es demasiado ancha indica high variance, si es demasiado angosta indica bias alto;
    \item \textbf{Curva de Critic loss}: un aumento abrupto significa que el learning rate es demasiado alto o que el objetivo es inestable;
    \item \textbf{Entropy}: si cae demasiado rápido indica que la exploration ha caído en determinism.
\end{itemize}

\begin{warningbox}{El tradeoff entre las métricas}
Reducir el critic loss no siempre es bueno: mientras el sesgo en la estimación de advantage sea grande, el Actor recibirá una señal incorrecta. Un entrenamiento estable requiere observar simultáneamente las tendencias del loss, del KL y de la policy entropy.
\end{warningbox}

\subsection{Resumen de la sección}

Las familias predominantes de Actor-Critic ponen énfasis distinto en el trust region, la entropy y el off-policy buffer; en la práctica es necesario observar varias curvas junto con early stopping o checkpoint rollback.

